\chapter{Schutz- und Robustheitsfunktionen}
\label{ch:schutz}
% Je Abschnitt: Anlass (ggf. Hardwarebefund des Laborpartners), Umsetzung
% (Skript, Parameter), Nachweis (test_config, Äquivalenz, MIL, Codegen, Hardware)

\section{Begrenzung des Bremsstroms (A2)}
% ctrl_limits.Iq_min (legacy -0,5 bitgleich, datasheet 0); Rückspeisung ins Labornetzteil
\offen{Text}

\section{Obergrenze des Drehzahlreglers}
% ctrl_limits.Iq_max
\offen{Text}

\section{Rücksetzen der Integratoren im Freemode}
% Stromstoß -31 A beim Freigeben (Windup), reset_I_cur_freemode, Drehzahlregler-Reset
\offen{Text, Telemetrie vorher/nachher}

\section{Sollwert-Totzone}
% Kriechen bei Poti 0, foc_speed_ref_deadband
\offen{Text}

\section{Eingangsqualifikation der Hall-Signale}
% 6 Samples, 500 ns (09_Hall_Qualifikation_2026-11-10.md)
\offen{Text}

\section{ADC-Abtastfenster}
% 300 ns für Ströme, 1 us für langsame Kanäle; Kalibrierblöcke in focLib
% (09_ADC_Abtastfenster_2026-11-10.md)
\offen{Text}

\section{Überspannungsabschaltung}
% 30 V, Hysterese 1 V, im Input Controller (0,2 s); Fehlerbit im FOCP-STATUS;
% Udc_max aus FOC Studio änderbar. Einschränkung: kein schneller Schutz (HW19)
% (09_Ueberspannung_2026-11-20.md)
\offen{Text}
